% part_2_model.tex — placeholder scaffold
% auto-mm-writing populates each section from the matching upstream artifact.

%--------------------------------- Data preprocessing
\section{Data Preprocessing}
% PLACEHOLDER — drawn from stage0_triage/data_recon.md and how anomalies
% were handled in stage2_solving/pipeline.py. Include a data-flow figure
% (fig-data-flow.pdf in img/) if available.
Data preprocessing description.

%--------------------------------- Model build (one section per sub-question)
\section{Task 1: <Title from problem restatement>}
\subsection{Model formulation}
% PLACEHOLDER — derived from stage1_modeling/model.md. Show notation,
% objective(s), constraint groups (each with a 1-paragraph English gloss),
% any non-trivial derivations.
Model description, formulas, and derivations.

\subsection{Algorithm}
% PLACEHOLDER — prose summary of stage2_solving/pipeline.py for this task.
% Each algorithmic component must answer "what problem feature does this address?"
% per integrity-rules.md Rule 5.
Algorithm description; cite stage2_solving/runs/<exp_id>/ for the implementation.

\subsection{Results}
% PLACEHOLDER — from stage2_solving/validation.md. Reference figures.
Results paragraph with hard numbers.

% (Repeat \section{Task N: ...} for each sub-question.)

%--------------------------------- Sensitivity
\section{Sensitivity Analysis}
% PLACEHOLDER — one subsection per `to-sweep` assumption from
% stage1_modeling/assumptions.md, populated from stage2_solving/sensitivity.md.
Sensitivity discussion.

%--------------------------------- Validation
\section{Model Validation}
% PLACEHOLDER — from stage2_solving/validation.md. Baseline comparison,
% small-instance exact Gap (if applicable), ablation findings.
Validation discussion.
